\section*{الفصل \ref{ch:g11:color-light-sources} --- اللون ومنابع الضوء}

\begin{solution}{exo:g11:color-light-sources:1}
(أ) أصفر؛ (ب) سماوي؛ (ج) أبيض. والمتمّم للأزرق: الأصفر
(أزرق $+$ أصفر $=$ أزرق $+$ أحمر $+$ أخضر $=$ أبيض).
\end{solution}

\begin{solution}{exo:g11:color-light-sources:2}
الأصفر: R وG. والأرجواني: R وB. والأبيض: R وG وB. والأسود: لا شيء. والبرتقالي: R
بكامل \omterm{def:g10:energy-conservation:power}{الاستطاعة} وG \omterm{def:g10:energy-conservation:power}{باستطاعة} مخفَّضة.
\end{solution}

\begin{solution}{exo:g11:color-light-sources:3}
يمتصّ الأرجواني الحزمة الخضراء: فيخرج $(R, B)$، أي ضوء أرجواني. ثم
يمتصّ \omterm{def:g11:color-light-sources:subtractive}{المرشّح} السماوي الحزمة الحمراء: فلا يبقى إلا $(B)$ --- أي أزرق.
\end{solution}

\begin{solution}{exo:g11:color-light-sources:4}
الشريط الأصفر: \omterm{def:g11:color-light-sources:objectcolor}{ينثر} R وG، ويمتصّ B. (أ) وتحت ضوء أحمر
\omterm{def:g11:color-light-sources:objectcolor}{ينثر} R: فيبدو أحمر. (ب) وتحت ضوء أزرق يمتصّ كل شيء: فيبدو أسود.
\end{solution}

\begin{solution}{exo:g11:color-light-sources:5}
$T = \dfrac{\qty{2.90e-3}{m.K}}{\qty{5.80e-7}{m}} = \qty{5000}{K}$.
\end{solution}

\begin{solution}{exo:g11:color-light-sources:6}
\emph{(أ)} $\lambda_{\max} = \qty{2.90e-3}{m.K}/\qty{1800}{K}
= \qty{1.61e-6}{m} = \qty{1610}{nm}$: أي \omterm{def:g10:light-spectra:wavelength}{تحت أحمر}.

\emph{(ب)} \omterm{def:g10:light-spectra:continuous}{الطيف مستمر}: فذيله المرئي صغير لكنه ليس
معدومًا، والطرف الأزرق يخبو أولًا --- فما يبقى مرئيًّا معظمه
أحمر برتقالي، ومن هنا لون اللهب.
\end{solution}

\begin{solution}{exo:g11:color-light-sources:7}
السترة (تنثر R): (أ) حمراء؛ (ب) حمراء؛ (ج) سوداء. والقفازان (ينثران G و
B، ويمتصّان R): (أ) سماويان؛ (ب) أسودان؛ (ج) أخضران.
\end{solution}

\begin{solution}{exo:g11:color-light-sources:8}
\emph{(أ)} \omterm{def:g10:light-spectra:emission}{طيف إصدار خطي}. ولا يقدر أي جسم أن \omterm{def:g11:color-light-sources:objectcolor}{ينثر} إلا
$\qty{589}{nm}$ أو لا شيء: فالشارع درجات من الأصفر البرتقالي و
الأسود.

\emph{(ب)} السيارة الزرقاء لا تنثر إلا الحزمة الزرقاء. وضوء النهار يحويها:
فتبدو زرقاء. أما خط الصوديوم فلا يحمل أزرق: فيُمتصّ كل شيء،
وتبدو السيارة سوداء.
\end{solution}

\begin{solution}{exo:g11:color-light-sources:9}
\emph{(أ)} تنتشر الحزمة البيضاء في قوس قزح؛ أما حزمة \omterm{def:g11:color-light-sources:lamps}{الليزر} فتنحرف
لكنها تبقى حزمة واحدة بالأخضر نفسه: فهي أحادية اللون
--- أي طول \omterm{def:g10:signals-and-waves:wave}{موجة} واحد، فلا شيء يرتّبه الموشور.

\emph{(ب)} ينفذ \omterm{def:g11:color-light-sources:subtractive}{المرشّح} الأحمر الحزمة الحمراء وحدها ويمتصّ
$\qty{532}{nm}$: فلا يكاد يخرج شيء؛ وتختفي البقعة.
\end{solution}

\begin{solution}{exo:g11:color-light-sources:10}
\omterm{def:g10:light-spectra:white}{الطيف}: ذروة زرقاء ضيقة قرب $\qty{450}{nm}$ مع حزمة \omterm{def:g11:color-light-sources:led}{الفوسفور}
الصفراء العريضة (نحو $500$--$\qty{700}{nm}$). والأزرق والأصفر
متممان، فيقرأ الثلاثي المخروطي أبيض. لكن الحزمة تخبو نحو
الأحمر الداكن، فلا يتلقى الجسم الأحمر الداكن إلا القليل مما يقدر أن ينثره: فيبدو
أبهت منه في ضوء النهار، الذي يغذّيه طيفه المستمر تغذية كاملة.
\end{solution}

\begin{solution}{exo:g11:color-light-sources:11}
(أ) $\lambda_{\max} = \qty{2.90e-3}{m.K}/\qty{1800}{K}
= \qty{1610}{nm}$، أي \omterm{def:g10:light-spectra:wavelength}{تحت أحمر}؛
(ب) $\qty{2.90e-3}{m.K}/\qty{6500}{K} = \qty{446}{nm}$، أي أزرق بنفسجي.
ويسمّى ضوء الشمعة ``دافئًا'' --- وهو المنبع ذو درجة الحرارة
\emph{الأدنى}: فالمفردات تسير عكس سُلّم \omterm{def:g10:light-spectra:kelvin}{الكلفن}.
\end{solution}

\begin{solution}{exo:g11:color-light-sources:12}
\emph{(أ)} متماثلان لونيان. ويبرهن وجودهما على أن اللون \omterm{def:g11:color-light-sources:perception}{إدراك}
--- أي ثلاث استجابات مخروطية --- لا \omterm{def:g10:light-spectra:white}{الطيف} نفسه: إذ ترتسم
أطياف لا تحصى على الثلاثي نفسه.

\emph{(ب)} موشور أو مِطياف: فخط واحد عند $\qty{589}{nm}$ في مقابل
حزمتين عند $540$ و$\qty{610}{nm}$. أو \omterm{def:g11:color-light-sources:subtractive}{مرشّح} أحمر: فهو يمتصّ
$\qty{589}{nm}$، فتظلم بقعة الصوديوم، بينما يبقى بكسل الشاشة
$\qty{610}{nm}$ --- فتصير رقعة الشاشة حمراء.
\end{solution}

\begin{solution}{exo:g11:color-light-sources:13}
\emph{(أ)} $\lambda_{\max} = \qty{2.90e-3}{m.K}/\qty{3200}{K} =
\qty{906}{nm}$. \omterm{def:g11:color-light-sources:colortemp}{ودرجة حرارة اللون} $\qty{3200}{K} < \qty{6500}{K}$:
أي أدفأ (أميل إلى البرتقالي) من ضوء النهار.

\emph{(ب)} يمتصّ \omterm{def:g11:color-light-sources:subtractive}{المرشّح} جزءًا من الفائض الأحمر البرتقالي ليعيد موازنة
\omterm{def:g10:light-spectra:white}{الطيف} نحو الأزرق؛ وبما أن \omterm{def:g11:color-light-sources:subtractive}{المرشّح} لا ينزع إلا الضوء، تهبط
\omterm{def:g10:energy-conservation:power}{الاستطاعة} الكلية على المسرح.

\emph{(ج)} \omterm{def:g11:color-light-sources:subtractive}{المرشّح} سلبي: فهو يقدر أن يمتصّ عند كل طول \omterm{def:g10:signals-and-waves:wave}{موجة} لكنه
لا يصدر أبدًا، فلا \omterm{def:g11:color-light-sources:subtractive}{مرشّح} يقدر أن يرفع شدة الأزرق فوق ما ينتجه
الكشّاف أصلًا.
\end{solution}

\begin{solution}{exo:g11:color-light-sources:14}
\emph{(أ)} $\qty{2700}{K}$: ضوء دافئ غني بالأحمر --- فالحزمة الحمراء في اللحم
مغذّاة جيدًا، فيبدو اللحم طريًّا زاهيًا. و$\qty{6500}{K}$: أبيض
مزرقّ بارد --- \omterm{def:g11:color-light-sources:objectcolor}{فينثر} السمك الثلاثي كاملًا ويبدو أبيض
ناصعًا، كأنه مثلّج للتوّ.

\emph{(ب)} لم يغيّر أيٌّ منهما بضاعته: فعادات الامتصاص في اللحم
والسمك لم تُمسّ. إنما تغيرت الإضاءة، ومن ثم الضوء المنثور ---
\omterm{def:g11:color-light-sources:objectcolor}{فلون الجسم} لا يُعرَّف إلا تحت إضاءة
معطاة.
\end{solution}

\begin{solution}{exo:g11:color-light-sources:15}
\emph{(أ)} الصوديوم: فوهج المدينة كله يقع في خط واحد عند
$\qty{589}{nm}$، وينزعه \omterm{def:g11:color-light-sources:subtractive}{مرشّح} ضيق واحد من صور
المقاريب. أما \omterm{def:g10:light-spectra:white}{طيف} LED العريض فلا يمكن ترشيحه دون التخلص
من ضوء النجوم أيضًا.

\emph{(ب)} يفضّل السائقون والمشاة LED: فتحت طيفه الكامل
تحتفظ الأجسام بألوان قابلة للتمييز (فالمعطف الأحمر يبقى أحمر)،
بينما يجعل ضوء الصوديوم كل شيء أصفر برتقاليًّا أو أسود.

\emph{(ج)} ثنائيات LED بيضاء دافئة، نحو $\qty{2700}{K}$: أي أداء لوني
مقبول للشارع، \omterm{def:g10:light-spectra:white}{وطيف} فقير في الطرف الأزرق، وهو الجزء
الأشدّ تلويثًا لسماء الليل.
\end{solution}

\begin{solution}{pb:g11:color-light-sources:1}
\textbf{1.} R $+$ G $=$ أصفر، وR $+$ B $=$ أرجواني، وG $+$ B $=$ سماوي؛
والتراكب الثلاثي أبيض.

\textbf{2.} البرتقالي: الأحمر بكامل \omterm{def:g10:energy-conservation:power}{الاستطاعة}، والأخضر \omterm{def:g10:energy-conservation:power}{باستطاعة} جزئية، والأزرق
مطفأ. والوردي الفاتح: الثلاثة مضاءة (أساس أبيض) مع غلبة طفيفة للأحمر.

\textbf{3.} الظلّ الملقى من الكشّاف الأحمر لا يتلقى إلا الكشّاف
الأخضر: فهو ظلّ أخضر؛ وبالتناظر يكون الظلّ الآخر أحمر.

\textbf{4.} الدهان الأزرق الداكن لا \omterm{def:g11:color-light-sources:objectcolor}{ينثر} إلا B. فتحت الكشّاف الأحمر وحده:
أسود. وتحت الثلاثة معًا: \omterm{def:g11:color-light-sources:objectcolor}{ينثر} الحزمة الزرقاء --- أي أزرق.

\textbf{5.} تختزل العين أي \omterm{def:g10:light-spectra:white}{طيف} إلى ثلاث استجابات مخروطية؛
وتوافق الثلاثي يكفي لتوافق اللون. وثلاثة أنواع من \omterm{def:g11:color-light-sources:cones}{المخاريط} تثبّت
عدد الكشّافات.

\textbf{6.} أصفر $=$ أبيض $-$ أزرق: فالصبغة تمتصّ الحزمة الزرقاء.

\textbf{7.} الكشّاف الأحمر: \omterm{def:g11:color-light-sources:objectcolor}{ينثر} R --- أي أحمر. والكشّاف الأخضر: \omterm{def:g11:color-light-sources:objectcolor}{ينثر} G ---
أي أخضر. والكشّاف الأزرق: يمتصّ كل شيء --- أي أسود.

\textbf{8.} أضِئ المشهد بالكشّاف الأزرق وحده؛ فالنقرة تحوّل
من الأبيض $\to$ الأزرق. فالزيّ الأصفر \omterm{def:g11:color-light-sources:objectcolor}{ينثر} R وG لكنه يمتصّ B:
فهو ساطع تحت الأبيض، أسود تحت الأزرق.

\textbf{9.} \omterm{def:g11:color-light-sources:subtractive}{المرشّح} الأرجواني: $(R, G, B) \to (R, B)$. \omterm{def:g11:color-light-sources:objectcolor}{وينثر} الزيّ
R وG مما يتلقاه: \omterm{def:g11:color-light-sources:objectcolor}{فينثر} R، ويمتصّ B --- أي زيّ أحمر.

\textbf{10.} \omterm{def:g11:color-light-sources:subtractive}{المرشّح} الأحمر يبقي حزمة من ثلاث ويرمي ثلثي
\omterm{def:g10:energy-conservation:power}{استطاعة} الكشّاف الأبيض؛ أما كل \omterm{def:g11:color-light-sources:subtractive}{مرشّح} C أو M أو Y فينزع حزمة واحدة، ويبقي
الثلثين، وتكديس اثنين منها يعطي مع ذلك أي لون أساسي.

\textbf{11.} \omterm{def:g10:light-spectra:emission}{طيف إصدار خطي}. ويصير الشارع عالمًا \omterm{def:g10:light-spectra:white}{أحادي
اللون}: فكل سطح أصفر برتقالي، أخفت أو أسطع، أو أسود.

\textbf{12.} الجدار الأبيض: أصفر برتقالي (\omterm{def:g11:color-light-sources:objectcolor}{ينثر} الخط). والعلامة
الصفراء: أصفر برتقالي ساطع. والسيارة الحمراء: داكنة --- فالخط $\qty{589}{nm}$
يقع خارج الحزمة التي تنثرها. والسيارة الزرقاء: سوداء.

\textbf{13.} \omterm{def:g11:color-light-sources:subtractive}{مرشّح} ضيق واحد مركزه $\qty{589}{nm}$ ينزع
وهج المدينة كله من صورة دون التضحية إلا بالنزر اليسير من
أطياف النجوم العريضة.

\textbf{14.} $\lambda_{\max} = \qty{2.90e-3}{m.K}/\qty{2700}{K} =
\qty{1.07e-6}{m} = \qty{1070}{nm}$: فالذروة، ومعها معظم \omterm{def:g10:energy-conservation:power}{الاستطاعة}،
تقع في \omterm{def:g10:light-spectra:wavelength}{تحت الأحمر} --- أي حرارة لا ضوءًا.

\textbf{15.} الصوديوم: $0.30 \times \qty{100}{W} = \qty{30}{W}$ من
الضوء المرئي؛ والتنغستن: $0.05 \times \qty{100}{W} = \qty{5}{W}$. أي ضوء أكثر ست
مرات بالفاتورة نفسها.

\textbf{16.} يحوّل \omterm{def:g11:color-light-sources:led}{الفوسفور} جزءًا من الأزرق إلى حزمة صفراء
عريضة؛ والأزرق والأصفر متممان، فيكمل الأزرق $+$ الأصفر
الثلاثي: أي أبيض.

\textbf{17.} سنّ أزرق ضيق مع حدبة صفراء عريضة --- في مقابل
الاتصال المتساوي لضوء النهار. والطرف الأحمر الداكن (مع انخفاض بين الأزرق و
الأصفر) ناقص التمثيل.

\textbf{18.} تُعرَض الأجسام الحمراء الداكنة باهتة مائلة إلى البنّي. ومن هنا
ثنائيات الجزّار الدافئة $\qty{2700}{K}$، التي يبقي محتواها الأحمر المعزَّز
اللحم زاهيًا.

\textbf{19.} لا شيء بداخله عند $\qty{2700}{K}$: فالرقاقة تشتغل قرب
درجة حرارة الغرفة. واللافتة لا تَعِد إلا بمثل \emph{صبغة} جسم
متوهّج عند $\qty{2700}{K}$، تكون ذروته
$\qty{2.90e-3}{m.K}/\qty{2700}{K} = \qty{1.07e-6}{m} = \qty{1070}{nm}$.

\textbf{20.} غرفة النوم: نحو $\qty{2700}{K}$ --- أي ضوء دافئ فقير بالأزرق
للراحة. والجواهري: $5500$--$\qty{6500}{K}$ --- أي أبيض شبيه بضوء النهار، حتى
تُقرأ البيض بيضًا وتردّ الأحجار كل حزمة. والقانون: اللون
المدرَك $=$ \omterm{def:g10:light-spectra:white}{طيف} المنبع $\times$ امتصاص الجسم، مقروءًا عبر
ثلاثة مخاريط.
\end{solution}
